\documentclass[border=0pt]{standalone}
% The shared preamble of every docs diagram (docs/diagrams/<name>.tex).
% dvisvgm's driver first, so patterns and opacity survive into the SVG.
\def\pgfsysdriver{pgfsys-dvisvgm.def}
\usepackage[siunitx,RPvoltages]{circuitikz}
\usepackage{xcolor}
% Latin Modern: outline fonts for every glyph a label uses, µ included.
\usepackage[T1]{fontenc}
\usepackage{lmodern}
\usetikzlibrary{arrows.meta,backgrounds,calc,patterns,positioning}
\tikzset{
  % A white card, so diagrams read in mdBook's dark themes too.
  every picture/.append style={show background rectangle,
    background rectangle/.style={fill=white,inner sep=4pt}},
  % Layout geometry: layers told apart by fill and hatching, never colour alone.
  layer a/.style={draw=black,line width=0.6pt,fill=gray!20},
  layer b/.style={draw=black,line width=0.6pt,fill=gray!45},
  layer c/.style={draw=black,line width=0.6pt,fill=white,postaction={pattern=north east lines}},
  layer d/.style={draw=black,line width=0.6pt,fill=white,postaction={pattern=crosshatch dots}},
  marker/.style={draw=black,dashed,line width=0.6pt},
  % A measured distance and its label.
  dim/.style={{Latex[length=2mm]}-{Latex[length=2mm]},line width=0.5pt},
  dim label/.style={font=\small,fill=white,inner sep=1pt},
  % Where a rule fails: thick, dashed and labelled, not only red.
  violation/.style={draw=red!75!black,line width=1.4pt,dash pattern=on 3pt off 2pt},
  violation label/.style={font=\small\bfseries,text=red!75!black},
  caption/.style={font=\small},
}

\begin{document}
\begin{circuitikz}
  % 1600 um die, 0.004 cm per um: 6.4 cm. metal1 everywhere but the upper-right quadrant,
  % which holds one 200 um block.
  \draw[layer a] (0,0) -- (6.4,0) -- (6.4,3.2) -- (3.2,3.2) -- (3.2,6.4) -- (0,6.4) -- cycle;
  \draw[layer a] (4,4) rectangle (4.8,4.8);
  \node[caption] at (1.6,1.6) {metal1};
  \draw[marker] (0,0) rectangle (6.4,6.4);
  \node[caption,anchor=north west] at (0,-0.05) {prboundary};
  % window (0,0) and the next one along x, step 400 um
  \draw[marker,line width=1pt] (0,0) rectangle (3.2,3.2);
  \draw[marker,line width=1pt] (1.6,0.1) rectangle (4.8,3.3);
  % the failing window, (800 um, 800 um) to (1600 um, 1600 um)
  \draw[violation] (3.2,3.2) rectangle (6.4,6.4);
  \node[violation label,anchor=south] at (4.8,6.5) {M1Fil.h: $d = 200^2/800^2 = 0.0625 < 0.25$};
  \draw[dim] (0,-0.75) -- node[dim label,below=1pt] {window \SI{800}{\micro\metre}} (3.2,-0.75);
  \draw[dim] (0,-1.55) -- node[dim label,below=1pt] {step \SI{400}{\micro\metre}} (1.6,-1.55);
  \draw[dim] (6.9,0) -- node[dim label,right=1pt] {\SI{1600}{\micro\metre}} (6.9,6.4);
  \draw[dim] (4,3.8) -- node[dim label,below=1pt] {\SI{200}{\micro\metre}} (4.8,3.8);
\end{circuitikz}
\end{document}
